% Préambule commun des schémas TikZ — charte Liouaï
% Chaque schéma commence par % Préambule commun des schémas TikZ — charte Liouaï
% Chaque schéma commence par % Préambule commun des schémas TikZ — charte Liouaï
% Chaque schéma commence par \input{_preambule} puis \begin{document}.
\documentclass[tikz,border=6pt]{standalone}
\usepackage[T1]{fontenc}
\usepackage[utf8]{inputenc}
\usepackage[scaled=0.95]{helvet}
\renewcommand{\familydefault}{\sfdefault}
\usepackage{amsmath,amssymb}
\usepackage[european]{circuitikz}
\usetikzlibrary{arrows.meta,positioning,calc,angles,quotes,decorations.pathreplacing,backgrounds}

% Couleurs du logo Liouaï Digital Education
\definecolor{navy}{HTML}{243743}
\definecolor{vert}{HTML}{28B78D}
\definecolor{vertfonce}{HTML}{167A5C}
\definecolor{vertclair}{HTML}{EAF7F2}
\definecolor{gris}{HTML}{A1A9AE}
\definecolor{gristexte}{HTML}{5C6870}
\definecolor{fond}{HTML}{F3F5F6}
\definecolor{jaune}{HTML}{E5C800}
\definecolor{jauneclair}{HTML}{FFFBD6}

\tikzset{
  every picture/.style={line cap=round, line join=round, color=navy, text=navy},
  fleche/.style={-{Stealth[length=2.6mm,width=2mm]}, line width=1.1pt},
  vecteur/.style={fleche, line width=1.4pt},
  titre/.style={font=\bfseries\small, text=navy},
  note/.style={font=\footnotesize, text=gristexte},
  boite/.style={draw=navy, line width=0.9pt, rounded corners=5pt, fill=fond, align=center, inner sep=6pt},
}
 puis \begin{document}.
\documentclass[tikz,border=6pt]{standalone}
\usepackage[T1]{fontenc}
\usepackage[utf8]{inputenc}
\usepackage[scaled=0.95]{helvet}
\renewcommand{\familydefault}{\sfdefault}
\usepackage{amsmath,amssymb}
\usepackage[european]{circuitikz}
\usetikzlibrary{arrows.meta,positioning,calc,angles,quotes,decorations.pathreplacing,backgrounds}

% Couleurs du logo Liouaï Digital Education
\definecolor{navy}{HTML}{243743}
\definecolor{vert}{HTML}{28B78D}
\definecolor{vertfonce}{HTML}{167A5C}
\definecolor{vertclair}{HTML}{EAF7F2}
\definecolor{gris}{HTML}{A1A9AE}
\definecolor{gristexte}{HTML}{5C6870}
\definecolor{fond}{HTML}{F3F5F6}
\definecolor{jaune}{HTML}{E5C800}
\definecolor{jauneclair}{HTML}{FFFBD6}

\tikzset{
  every picture/.style={line cap=round, line join=round, color=navy, text=navy},
  fleche/.style={-{Stealth[length=2.6mm,width=2mm]}, line width=1.1pt},
  vecteur/.style={fleche, line width=1.4pt},
  titre/.style={font=\bfseries\small, text=navy},
  note/.style={font=\footnotesize, text=gristexte},
  boite/.style={draw=navy, line width=0.9pt, rounded corners=5pt, fill=fond, align=center, inner sep=6pt},
}
 puis \begin{document}.
\documentclass[tikz,border=6pt]{standalone}
\usepackage[T1]{fontenc}
\usepackage[utf8]{inputenc}
\usepackage[scaled=0.95]{helvet}
\renewcommand{\familydefault}{\sfdefault}
\usepackage{amsmath,amssymb}
\usepackage[european]{circuitikz}
\usetikzlibrary{arrows.meta,positioning,calc,angles,quotes,decorations.pathreplacing,backgrounds}

% Couleurs du logo Liouaï Digital Education
\definecolor{navy}{HTML}{243743}
\definecolor{vert}{HTML}{28B78D}
\definecolor{vertfonce}{HTML}{167A5C}
\definecolor{vertclair}{HTML}{EAF7F2}
\definecolor{gris}{HTML}{A1A9AE}
\definecolor{gristexte}{HTML}{5C6870}
\definecolor{fond}{HTML}{F3F5F6}
\definecolor{jaune}{HTML}{E5C800}
\definecolor{jauneclair}{HTML}{FFFBD6}

\tikzset{
  every picture/.style={line cap=round, line join=round, color=navy, text=navy},
  fleche/.style={-{Stealth[length=2.6mm,width=2mm]}, line width=1.1pt},
  vecteur/.style={fleche, line width=1.4pt},
  titre/.style={font=\bfseries\small, text=navy},
  note/.style={font=\footnotesize, text=gristexte},
  boite/.style={draw=navy, line width=0.9pt, rounded corners=5pt, fill=fond, align=center, inner sep=6pt},
}

\begin{document}
\begin{tikzpicture}[scale=0.9]
\draw[fleche, gris, line width=0.8pt] (0,0) -- (8.3,0) node[note, below left] {hauteur de lâcher $h$};
\draw[fleche, gris, line width=0.8pt] (0,0) -- (0,5.4) node[note, above right] {hauteur du rebond $r$};
\draw[vertfonce, line width=1.6pt] (0.3,0.45) -- (7.9,4.95) node[right, font=\small, text=vertfonce] {$r = a\,h + b$};
\foreach \x/\y in {1/1.35, 2/1.55, 3/2.55, 4/2.55, 5/3.6, 6/3.9, 7/4.8}{
  \pgfmathsetmacro{\yd}{0.45 + (\x-0.3)*4.5/7.6}
  \draw[jaune!85!navy, line width=1.2pt] (\x,\y) -- (\x,\yd);
  \fill[navy] (\x,\y) circle (2.4pt);
}
\node[font=\footnotesize, text=navy, anchor=west] at (4.1,1.75) {résidu $e_i = r_i - (a\,h_i + b)$};
\draw[gristexte, line width=0.5pt] (4.1,1.85) -- (4.0,2.4);
\node[boite, anchor=west, align=left, font=\small] at (9.2,2.6) {Moindres carrés : on choisit $a$ et $b$\\ qui rendent la somme\\[3pt]
$\displaystyle S(a,b) = \sum_i e_i^2$\\[3pt] la plus petite possible.\\[4pt]
\textcolor{gristexte}{Les résidus sont verticaux :}\\ \textcolor{gristexte}{on mesure l'erreur sur $r$.}};
\end{tikzpicture}
\end{document}
